% Magnetar graphical abstract for the poster — adapted from
% manuscript/figures/gertsenshtein_zeldovich_star.tex and recoloured to the
% modern Cambridge poster palette:
%   incoming gravitational wave -> Dark Blue
%   dipole magnetic field lines -> dark teal (Warm/Dark Blue)
%   neutron star                -> Dark Blue
%   magnetosphere / corona      -> Light Blue
%   outgoing photon / FRB       -> Cherry (the amplification colour)
\documentclass[border=3pt,tikz]{standalone}
\usepackage{bm}
\usepackage{tikz,pgfplots}
\usetikzlibrary{arrows.meta}
\tikzset{>=latex}
\pgfplotsset{compat=1.13}
\usetikzlibrary{decorations.markings,intersections,calc}
\usepackage{ifthen}
\usepackage[outline]{contour}  % xcolor already loaded by the standalone 'tikz' option

% ---- modern Cambridge palette ----
\definecolor{CamBlue}{HTML}{8EE8D8}
\definecolor{CamLightBlue}{HTML}{D1F9F1}
\definecolor{CamWarmBlue}{HTML}{00BDB6}
\definecolor{CamDarkBlue}{HTML}{133844}
\definecolor{CamCherry}{HTML}{CD3572}

\colorlet{Bcol}{CamWarmBlue}
\colorlet{BcolFL}{CamDarkBlue!78!CamWarmBlue}  % dipole field lines: dark teal
\colorlet{Corona}{CamBlue}                     % magnetosphere halo
\colorlet{GW}{CamDarkBlue}                      % incoming gravitational wave
\colorlet{Photon}{CamCherry}                    % outgoing photon / FRB
\contourlength{1.2pt}

\usetikzlibrary{fadings}
\begin{tikzfadingfrompicture}[name=halo]
  \shade[inner color=transparent!0,outer color=transparent!100] (0,0) circle (0.9);
\end{tikzfadingfrompicture}

\tikzstyle{gw}=[-{Latex[length=2mm,width=1.5mm]}, GW, thick,decorate, decoration={snake,amplitude=3,segment length=8,post length=8}, path fading=east]


\begin{document}

\begin{tikzpicture}[
    line cap=round,
    line join=round]
    \def\xmax{8}
    \def\ymax{2}
    \def\E{0.7}
    \def\N{13}
    \def\cone{27}
    \def\tilt{8}

    \fill[path fading=halo,Corona] (0,0) circle (2*\E); % Halo/corona

    \draw[thick, dotted, Corona] (0,0) ellipse ({2*\E} and {1.7*\E});
    \fill[CamWarmBlue, opacity=0.075] (0,0) ellipse ({2*\E} and {1.7*\E});

    \clip (-\xmax,-\ymax) rectangle (\xmax,\ymax);
        \foreach \i [evaluate={\r=\i*\xmax;}] in {0.15, 0.3, 0.6, 0.99}{
            \draw[BcolFL, thick, domain=0:360, smooth, samples=50, variable=\t] plot ({\r*sin(\t)*sin(\t)*sin(\t)},{\r*sin(\t)*sin(\t)*cos(\t)});
        }
        \foreach \i/\ang [evaluate={\r=\i*\xmax;}] in {0.15/5, 0.3/2.6, 0.6/1.5, 0.99/0.8}{
            \draw[BcolFL, -{Latex[length=2mm,width=1.5mm]}, thin, domain=0:90+\ang, smooth, samples=20, variable=\t] plot ({\r*sin(\t)*sin(\t)*sin(\t)},{\r*sin(\t)*sin(\t)*cos(\t)});
            \draw[BcolFL, -{Latex[length=2mm,width=1.5mm]}, thin, domain=360:270-\ang, smooth, samples=20, variable=\t] plot ({\r*sin(\t)*sin(\t)*sin(\t)},{\r*sin(\t)*sin(\t)*cos(\t)});
        }
        \foreach \i/\ang [evaluate={\r=\i*\xmax;}] in {1.6/30, 3.0/20, 6.0/15, 20.0/10, 70.0/5}{
            \draw[BcolFL, thick, domain=0:\ang, smooth, samples=10, variable=\t] plot ({\r*sin(\t)*sin(\t)*sin(\t)},{\r*sin(\t)*sin(\t)*cos(\t)});
            \draw[BcolFL, thick, domain=180-\ang:180, smooth, samples=10, variable=\t] plot ({\r*sin(\t)*sin(\t)*sin(\t)},{\r*sin(\t)*sin(\t)*cos(\t)});
            \draw[BcolFL, thick, domain=180:180+\ang, smooth, samples=10, variable=\t] plot ({\r*sin(\t)*sin(\t)*sin(\t)},{\r*sin(\t)*sin(\t)*cos(\t)});
            \draw[BcolFL, thick, domain=360-\ang:360, smooth, samples=10, variable=\t] plot ({\r*sin(\t)*sin(\t)*sin(\t)},{\r*sin(\t)*sin(\t)*cos(\t)});
        }
        \foreach \i/\ang [evaluate={\r=\i*\xmax;}] in {0.3/58, 0.6/36.2, 1.6/20.83, 6.0/10.55}{
            \draw[BcolFL,-{Latex[length=2mm,width=1.5mm]}, thin, domain=0:\ang, smooth, samples=10, variable=\t] plot ({\r*sin(\t)*sin(\t)*sin(\t)},{\r*sin(\t)*sin(\t)*cos(\t)});
            \draw[BcolFL,-{Latex[length=2mm,width=1.5mm]}, thin, domain=90:180-\ang, smooth, samples=10, variable=\t] plot ({\r*sin(\t)*sin(\t)*sin(\t)},{\r*sin(\t)*sin(\t)*cos(\t)});
            \draw[BcolFL,-{Latex[length=2mm,width=1.5mm]}, thin, domain=270:180+\ang, smooth, samples=10, variable=\t] plot ({\r*sin(\t)*sin(\t)*sin(\t)},{\r*sin(\t)*sin(\t)*cos(\t)});
            \draw[BcolFL,-{Latex[length=2mm,width=1.5mm]}, thin, domain=360:360-\ang, smooth, samples=10, variable=\t] plot ({\r*sin(\t)*sin(\t)*sin(\t)},{\r*sin(\t)*sin(\t)*cos(\t)});
        }
        \draw[BcolFL, -{Latex[length=2mm,width=1.5mm]}, thin] (0,0) -- (0,0.803*\ymax);
        \draw[BcolFL, -{Latex[length=2mm,width=1.5mm]}, thin] (0,-\ymax) -- (0,-0.803*\ymax);

        \draw[BcolFL, thick] (0,0) -- (0,\ymax);
        \draw[BcolFL, thick] (0,0) -- (0,-\ymax);

    \draw[line width=1.5pt, white] (\tilt+90:1.1*\ymax) -- (\tilt-90:1.1*\ymax);

    \draw[very thin,Corona,fill=CamDarkBlue] (0,0) circle (\E); % Star
    \draw[thick, draw=Corona,ball color=CamDarkBlue,fill opacity=0.15] (0,0) circle (\E);
    \draw[dashed, thick, CamDarkBlue] (\tilt+90:1.1*\ymax) -- (\tilt-90:1.1*\ymax);

    \node[CamDarkBlue] at (0.203*\xmax,0.87*\ymax) {\contour{white}{Magnetosphere}};
    \node[CamDarkBlue] at (0.24*\xmax,-0.315*\ymax) {\contour{white}{Neutron star}};

    \foreach \i[evaluate={\y=\i*0.2*\ymax;}] in {-1,0,1}{
        \draw[gw] (-0.9*\xmax,\y) -- (-2*\E,\y);
    }
    \node[GW] at (-0.65*\xmax,0.4*\ymax) {\contour{white}{Gravitational waves}};

    \coordinate (EM) at (2*\E,0);
    \foreach \i [evaluate={\r={\i*0.57*\xmax/\N};}] in {1,...,\N}{
        \draw[line width=2.5pt, white] ($(EM)+(+\cone:{\r} and {0.5*\r})$) arc(+\cone:-\cone:{\r} and {0.5*\r});
        \draw[line width=1.75pt, Photon, opacity=0.12] ($(EM)+(+\cone:{\r} and {0.5*\r})$) arc(+\cone:-\cone:{\r} and {0.5*\r});
        \draw[thick, Photon] ($(EM)+(+\cone:{\r} and {0.5*\r})$) arc(+\cone:-\cone:{\r} and {0.5*\r});
    }

    \node[Photon,align=center] at (0.81*\xmax,0.55*\ymax) {\contour{white}{Photon} \\ \contour{white}{(Fast Radio Burst)}};
\end{tikzpicture}

\end{document}
